\section{Задания к главе}

\subsection*{Уровень 1: Обязательный}
\begin{enumerate}
    \item \textbf{Калькулятор полета.} Напишите программу, которая запрашивает у пользователя время полета (в минутах) и среднюю скорость (в м/с), а затем выводит пройденное расстояние в метрах и километрах.
    \item \textbf{Чек-лист.} Создайте список (\texttt{list}) из 5 пунктов проверки перед вылетом (например, <<Батарея заряжена>>, <<Винты целы>>). Используя цикл \texttt{for}, выведите каждый пункт с номером.
\end{enumerate}

\subsection*{Уровень 2: Усложненный}
\begin{enumerate}
    \item \textbf{Умная батарея.} Напишите функцию \texttt{check\_battery(voltage)}, которая принимает напряжение и возвращает статус:
    \begin{itemize}
        \item $> 4.0$ В: <<Отличный заряд>>
        \item $3.7 - 4.0$ В: <<Средний заряд>>
        \item $< 3.7$ В: <<Посадка немедленно!>>
    \end{itemize}
    Интегрируйте эту функцию в скрипт подключения к дрону.
\end{enumerate}

\subsection*{Уровень 3: Творческий (Мини-проект)}
\textbf{Миссия <<Квадрат>>.} 
Используя команды \texttt{drone.go\_to\_local\_point(x, y, z, yaw)}, напишите программу, которая заставляет дрон пролететь по траектории квадрата со стороной 1 метр на высоте 1.5 метра.
\textit{Подсказка:} Вам понадобятся координаты точек: $(0,0) \to (1,0) \to (1,1) \to (0,1) \to (0,0)$. Не забывайте про \texttt{time.sleep()} между командами!
